% theme: rotation_and_diurnal_motion
\MajorQuestion{自転と日周運動}
\SetRowSpaceQanda

\Instruction{次の問いに答えなさい。}
\QQ{太陽が真南を通過することを\NamedBlank{culmination}という。}{南中}
\QQ{天体が、中心を通る軸を中心に自ら回転することを\NamedBlank{rotation}という。}{自転}
\QQ{地球の北極と南極を結ぶ、回転の中心となる軸を\NamedBlank{axis}という。}{地軸}
\QQ{星などの天体がはりついているように見える、見かけ上の球面を\NamedBlank{celestial_sphere}という。}{天球}
\QQ{地球は1日に360$^\circ$回転するので、1時間には約何度回転するか。}{15$^\circ$}

\Instruction{\strut}
\QQ{地球の自転によって、天体が1日に1回動くように見える運動を\NamedBlank{diurnal}という。}{日周運動}
\QQ{地球は、北極側から見てどちらの向きに回っているか。}{西から東}
\QQ{\RefBlank{culmination}したときの太陽の高度を何というか。}{南中高度}
\QQ{観測者の真上の点を何というか。}{天頂}
\QQ{北の空で、ほとんど動かず中心付近に見える星を何というか。}{北極星}

\Instruction{\strut}
\QQ{太陽や星は、全体としてどちらの向きに動くように見えるか。}{東から西}
\QQ{自分がいる地点での方向を表すものを\NamedBlank{direction}という。}{方位}
\QQ{\RefBlank{axis}を北側へ伸ばしたとき、\RefBlank{celestial_sphere}と交わる点を\NamedBlank{celestial_north}という。}{天の北極}
\QQ{太陽が真南を通過する時刻を何というか。}{正午}
\QQ{複数の星が決まった形に集まって見えるものを何というか。}{星座}

\Instruction{\strut}
\QQ{北の空の星は、\RefBlank{celestial_north}付近を中心に、どちら回りに動くように見えるか。}{反時計回り}
\QQ{観測者から見て、天頂と南北を結ぶ線を\NamedBlank{meridian}という。}{子午線}
\QQ{空と地面の境目として見える線を何というか。}{地平線}
\QQ{\RefBlank{axis}を南側へ伸ばしたとき、\RefBlank{celestial_sphere}と交わる点を何というか。}{天の南極}
\QQ{太陽が東の空へ出て、昼が始まるころを何というか。}{日の出}
